\documentclass[11pt,a4paper]{article}

\usepackage[margin=24mm]{geometry}
\usepackage{amsmath,amssymb,bm,booktabs,tabularx,array}
\usepackage{hyperref}
\usepackage{xcolor}

\newcommand{\vect}[1]{\bm{#1}}
\newcommand{\Tens}{\bm P}
\newcommand{\rank}{\operatorname{rank}}
\newcommand{\Null}{\mathcal N}
\newcolumntype{Y}{>{\raggedright\arraybackslash}X}

\title{What One and Two L/R/U/D Polarimeters Can Determine\\
from a Spin--1 Beam}
\author{Concise DPOLAR analysis note}
\date{}

\begin{document}
\maketitle

\begin{abstract}
A general spin--1 density matrix has eight real polarization degrees of freedom:
three vector components and five tensor components.  A single ideal L/R/U/D
polarimeter does not determine all eight.  With several polar scattering angles,
controlled normalization, and non-degenerate analyzing powers, one station can
determine at most six combinations.  Its exact cardinal-geometry nulls are the
longitudinal vector polarization $p_z$ and the tensor component $p_{xy}$.  A
second station can determine all eight only if known spin transport rotates both
null directions into combinations observed at the other station.  If particles
follow different paths, the transport is an ensemble of rotations rather than a
single rotation; this can reduce polarization, mix components, and introduce
additional nuisance parameters that lower identifiability.
\end{abstract}

\section{The short answer}

The spin--1 density matrix can be written in terms of
\begin{equation}
 \vect p=(p_x,p_y,p_z),\qquad
 \Tens=
 \begin{pmatrix}
 p_{xx}&p_{xy}&p_{xz}\\
 p_{xy}&p_{yy}&p_{yz}\\
 p_{xz}&p_{yz}&p_{zz}
 \end{pmatrix},\qquad
 p_{xx}+p_{yy}+p_{zz}=0.
\end{equation}
Thus there are $3+5=8$ real polarization degrees of freedom.

\begin{table}[htbp]
\centering
\caption{Maximum structural information under ideal calibration.  ``Multiple
$\theta$'' means at least two polar-angle groups whose analyzing powers have
different angular dependence.}
\small
\begin{tabularx}{\linewidth}{YcYY}
\toprule
Configuration & Maximum rank & Determinable & Remaining exact nulls\\
\midrule
One station, one $\theta$, known normalization
& $4/8$ & one common mode, two first harmonics, one diagonal harmonic
& $p_z$, $p_{xy}$, and two vector--tensor first-harmonic mixtures\\
One station, multiple $\theta$, shared normalization
& $6/8$ & $p_x,p_y,p_{xz},p_{yz},p_{xx}-p_{yy},p_{zz}$
& $p_z$ and $p_{xy}$\\
Two identical stations, aligned transport
& $6/8$ & the same six combinations with better precision
& the same two null directions\\
Two stations, generic non-longitudinal relative rotation
& up to $8/8$ & all vector and tensor components
& none, if the transported null spaces do not overlap\\
Two stations, relative rotation only about $z$
& up to $7/8$ & the tensor null may be lifted
& $p_z$ remains invariant and invisible\\
\bottomrule
\end{tabularx}
\end{table}

These are upper bounds.  Unknown luminosities, independent detector efficiencies,
backgrounds, or uncertain spin transport can reduce the effective rank.

\section{Why one station has only six independent combinations}

In the local beam frame, the repository cross-section convention is
\begin{align}
 \frac{\sigma(\theta,\phi)}{\sigma_0(\theta)}=1
 &+\frac32(p_x\sin\phi+p_y\cos\phi)A_y(\theta)\notag\\
 &+\frac23(p_{xz}\cos\phi-p_{yz}\sin\phi)A_{xz}(\theta)\notag\\
 &+\frac16\left[(p_{xx}-p_{yy})\cos2\phi
 -2p_{xy}\sin2\phi\right][A_{xx}(\theta)-A_{yy}(\theta)]\notag\\
 &+\frac12p_{zz}A_{zz}(\theta).
 \label{eq:response}
\end{align}
The four azimuthal harmonics have the interpretation
\begin{center}
\begin{tabular}{lll}
\toprule
Harmonic & Polarization combination & L/R/U/D information\\
\midrule
common & $p_{zz}A_{zz}$ & absolute or cross-angle rate\\
$\cos\phi$ & $\frac32p_yA_y+\frac23p_{xz}A_{xz}$ & $L-R$\\
$\sin\phi$ & $\frac32p_xA_y-\frac23p_{yz}A_{xz}$ & $U-D$\\
$\cos2\phi$ & $(p_{xx}-p_{yy})(A_{xx}-A_{yy})/6$ & $(L+R)-(U+D)$\\
$\sin2\phi$ & $-p_{xy}(A_{xx}-A_{yy})/3$ & zero at cardinal azimuths\\
\bottomrule
\end{tabular}
\end{center}

At one polar angle, the first harmonics contain inseparable vector--tensor pairs:
$(p_y,p_{xz})$ and $(p_x,p_{yz})$.  Several polar angles separate each pair if
the angular responses are non-collinear, schematically
\begin{equation}
 \det
 \begin{pmatrix}
 A_y(\theta_1)&A_{xz}(\theta_1)\\
 A_y(\theta_2)&A_{xz}(\theta_2)
 \end{pmatrix}\neq0.
 \label{eq:first-harmonic-condition}
\end{equation}
Similarly, $p_{zz}$ can be separated from one shared luminosity scale only if the
acceptance-averaged $A_{zz}$ responses differ between angular groups.  Independent
normalizations for every angle remove this information.

Two directions remain invisible even with many polar angles:
\begin{enumerate}
 \item $p_z$ does not occur in Eq.~\eqref{eq:response}; the reaction is sensitive
 only to transverse vector polarization.
 \item $p_{xy}$ multiplies $\sin2\phi$, which is exactly zero for
 $\phi=0^\circ,90^\circ,180^\circ,270^\circ$ and for symmetric finite openings
 centered on those azimuths.
\end{enumerate}

For the repository's current tensor-only calculation, vector polarization is
excluded from the fit.  The planned two-angle proton-singles response then has
rank $4/5$: it determines $p_{xx}-p_{yy}$, $p_{xz}$, $p_{yz}$, and $p_{zz}$,
while $p_{xy}$ is the exact null direction.  This is consistent with the general
$6/8$ result after the two transverse vector components are restored.

\section{What a second station adds}

Let the upstream density-matrix parameters be a common reference vector
$\vect x_0=(\vect p_0,\vect q_0)$, with three vector and five tensor coordinates.
Known spin transport to station $s$ acts as
\begin{equation}
 \vect p_s=R_s\vect p_0,\qquad
 \Tens_s=R_s\Tens_0R_s^{\mathsf T},\qquad
 \vect q_s=T^{(2)}(R_s)\vect q_0.
\end{equation}
If $H_s$ is the local station response, the combined response is
\begin{equation}
 J_{\rm combined}=
 \begin{pmatrix}
 H_1M_1\\ H_2M_2
 \end{pmatrix},\qquad
 M_s=R_s\oplus T^{(2)}(R_s),
\end{equation}
and its null space is
\begin{equation}
 \boxed{\Null_{\rm combined}
 =\ker(H_1M_1)\cap\ker(H_2M_2)}.
 \label{eq:null-intersection}
\end{equation}

Equation~\eqref{eq:null-intersection} is the practical criterion.
\begin{itemize}
 \item If the two stations have the same geometry and aligned spin transport,
 they measure the same six-dimensional subspace.  Counts increase, but rank does
 not.
 \item A relative rotation about $z$ can rotate the tensor $p_{xy}$ null into a
 visible diagonal tensor component.  However, it leaves $p_z$ unchanged, so the
 full density matrix still has at least one null direction.
 \item A generic rotation containing an $x$ or $y$ component can rotate $p_z$
 into $p_x$ or $p_y$ and can simultaneously rotate $p_{xy}$ into visible tensor
 components.  Full rank eight is possible if both transported nulls are
 complementary and the nuisance model does not absorb them.
\end{itemize}

For the immediate nominal pure-$p_{yy}$ experiment, the parameter space is much
smaller than eight dimensions.  One after-SRC station is sufficient for tensor
magnitude and stability monitoring and for the first-order tilt that generates
$p_{yz}$.  The orthogonal tilt generates $p_{xy}$ and is invisible at one
cardinal station.  A second station is justified when a known transport rotation
makes that missing tilt or another contamination mode observable.

\section{Particles following different paths}

The single-rotation model assumes every particle samples the same magnetic fields
and follows the same orbit.  A real beam has transverse position, angle, momentum,
and time coordinates $\xi$.  Each trajectory may therefore have its own rotation
$R(\xi)$.  The state measured at station 2 is an ensemble average,
\begin{align}
 \overline{\vect p}_2
 &=\int d\xi\,f(\xi)R(\xi)\vect p_1,\\
 \overline{\Tens}_2
 &=\int d\xi\,f(\xi)R(\xi)\Tens_1R^{\mathsf T}(\xi).
 \label{eq:ensemble-transport}
\end{align}
The averaged vector and tensor maps are generally not rotations: they need not be
orthogonal and can contract the polarization.  Coherent rotation and phase spread
are therefore different physical effects.

For illustration, if all trajectories rotate about the local beam $z$ axis but
their rotation angle has a Gaussian RMS spread $\sigma_\psi$, an azimuthal harmonic of order $m$
is reduced by
\begin{equation}
 \left\langle e^{im\psi}\right\rangle
 =e^{im\langle\psi\rangle}e^{-m^2\sigma_\psi^2/2}.
\end{equation}
Vector first harmonics are suppressed by $e^{-\sigma_\psi^2/2}$, while tensor
second harmonics are suppressed more rapidly, by $e^{-2\sigma_\psi^2}$.  This can
look like depolarization even though every individual particle undergoes a
perfectly coherent spin rotation.

Trajectory variation also affects the detector response itself:
\begin{itemize}
 \item beam offsets and angles shift the local detector azimuths and mix the
 $\sin\phi$, $\cos\phi$, $\sin2\phi$, and $\cos2\phi$ harmonics;
 \item momentum-dependent spin phase can correlate polarization with acceptance;
 \item apertures and losses can change the trajectory weights between stations;
 \item a trajectory-dependent efficiency can generate false L/R or U/D
 asymmetries.
\end{itemize}
The measured channel mean is therefore more accurately written as
\begin{equation}
 \mu_k=\int d\xi\,f(\xi)\,w_k(\xi)\,
 \sigma_k\!\left[R(\xi)\vect p_1,\,R(\xi)\Tens_1R^{\mathsf T}(\xi)\right],
\end{equation}
where $w_k(\xi)$ contains orbit acceptance and detector efficiency.

If $f(\xi)$ and $R(\xi)$ are unknown, their parameters become nuisances and the
effective polarization rank can decrease.  Two polarimeters alone cannot recover
an arbitrary path distribution.  A defensible two-station analysis should use:
\begin{enumerate}
 \item an orbit and spin transfer map over the measured beam phase space;
 \item beam-position, angle, and momentum diagnostics at both stations;
 \item trajectory-binned or event-by-event response calculations where possible;
 \item finite detector acceptance and calibrated relative efficiencies;
 \item ensemble-transport parameters included in the Fisher or likelihood fit.
\end{enumerate}

\section{Practical conclusion}

One well-calibrated, multi-angle L/R/U/D station does not reconstruct a general
spin--1 density matrix.  Under favorable analyzing-power and normalization
conditions it determines six of eight vector-plus-tensor degrees of freedom;
$p_z$ and $p_{xy}$ remain exact nulls.  For the present tensor-focused pure-$p_{yy}$
program, this is sufficient for the primary after-SRC magnitude and stability
measurement.

Two stations provide genuinely new information only when the known relative
spin transport makes their null spaces complementary.  A generic rotation with
a transverse axis can, in principle, give full eight-dimensional density-matrix
rank.  Non-uniform particle paths replace that rotation by an ensemble map and
can reduce polarization and identifiability.  The scientific case for a second
station should therefore be based on a phase-space-dependent spin-transfer model
and a quantified reduction of the combined null space, not only on increased
statistics.

\end{document}
